\begin{center}\begin{tikzpicture}[font=\small,>=stealth,line join=round]
\draw[gray!40,rounded corners=3pt] (0,0) rectangle (14.2,8.3);\node at (7.1,7.700000000000001) {两次轴对称，把整幅图横移两轴间距的两倍};
\begin{scope}[shift={(1.9181818181818182,3.822222222222222)},x=0.5454545454545454cm,y=0.12777777777777777cm]\draw[->,gray] (-1.5,0)--(9.5,0) node[right]{$x$};\draw[->,gray] (0,-17)--(0,19) node[above]{$y$};\draw[blue,thick,solid,domain=-1:9,samples=180] plot(\x,{14*cos(deg(pi*(\x-1)/3))});\fill[orange!85!black] (2,7) circle(2.2pt);\fill[orange!85!black] (8,7) circle(2.2pt);\foreach \a in {1,4}{\draw[violet,dashed] (\a,-16)--(\a,16);\node[above,violet] at (\a,16){$x=\a$};}\end{scope}\node[anchor=west,align=left,text width=5.0cm] at (8.4,4.0) {示例 $f=14\cos\frac{\pi(x-1)}3$\\[9pt]两轴 $x=1,4$\\[9pt]两次配对 $x\to2-x\to x+6$\\[9pt]$f(2)=f(8)=7$\\[9pt]周期为6（只保证一个）};
\node[text width=12.6cm,align=center] at (7.1,.65) {本例曲线取振幅14以满足已知值7；正文推导适用于所有满足题设的函数。};
\end{tikzpicture}\end{center}
